% BOOK-FILE-ID: f-d7df492a82ce
\subsection{二方向の階差}\label{節：二方向の階差}
非負整数$n,m$で数表の値$a_{n,m}$を表す.
\sectionref*{節：階差型}の和と差を,横と縦の二方向で使おう.

\paragraphbf{例：距離の二乗と階差}
原点から点$(n,m)$までの距離は
\[ a_{n,m}=\sqrt{n^2+m^2} \]
である.
根号を取り除くため$b_{n,m}=a_{n,m}^2$とすると,
一方向の差は$b_{n+1,m}-b_{n,m}=2n+1$となり,$m$によらない.
従って,さらにもう一方向へ差を取れば0になる.

\paragraphbf{0の差から一定の差へ}
この計算を広げ,境界の値
\[
  b_{n,0}=n^2,\qquad b_{0,m}=m^2
\]
と,一定の実数$c$について
\BookMathLayout{\[
    b_{n+1,m+1}-b_{n+1,m}-b_{n,m+1}+b_{n,m}=c
  \]}{\[
    \begin{aligned}
      b_{n+1,m+1}-b_{n+1,m}\\
      {}-b_{n,m+1}+b_{n,m}=c
    \end{aligned}
  \]}
とする.$b_{n+1,m+1}$は添字の和が小さい三つの値から決まるので,$n+m$についての帰納法で境界から一意に求まる.

横方向の差を$d_{n,m}=b_{n+1,m}-b_{n,m}$と置くと,
\[
  d_{n,m+1}-d_{n,m}=c,\qquad d_{n,0}=2n+1
\]
なので$d_{n,m}=2n+1+cm$である.これを横に足し合わせれば,
\[
  b_{n,m}=n^2+m^2+cnm
\]
となる.$c$が一定なので,\textbf{二方向の\bookterm{difference}{階差}を逆順に足し合わせられた}.

\paragraphbf{数表と図形}
長さが$1$のベクトルを単位ベクトルという.
二つの単位ベクトル$\boldsymbol u,\boldsymbol v$を同じ始点から描いたときの角を$\theta$とすると,
$c=2\cos\theta$の場合には
\[
  b_{n,m}=|n\boldsymbol u+m\boldsymbol v|^2
\]
と読める.数表なら$c$は任意の実数でよいが,この距離の二乗という解釈には$-2\le c\le2$が必要になる.

\paragraphbf{例：格子経路を数える}
一方,格子点$(0,0)$から$(n,m)$へ,右と上へ一歩ずつ進む道の数を$A_{n,m}$とする.
最後の一歩を分ければ,
\[
  A_{n,m}=A_{n-1,m}+A_{n,m-1}\quad(n,m\ge1)
\]
であり,境界では$A_{n,0}=A_{0,m}=1$である.
\BookFigLatticeRecurrence
下の数表は,上と左の値を足して作ったものである.
\[
  \begin{array}{c|rrrr}
    m\backslash n&0&1&2&3\\\hline
    0&1&1&1&1\\
    1&1&2&3&4\\
    2&1&3&6&10\\
    3&1&4&10&20
  \end{array}
\]
全$n+m$歩から上へ進む$m$歩の位置を選ぶので,
\[
  A_{n,m}=\binom{n+m}{m}
\]
である.最後の一歩の分類は,二項係数の加法公式でもある.
境界と帰納法による確認は\bookproblemref{practice-connections}{3}{接続演習問3}へ.

\paragraphbf{二つの例の共通点}
このように二つの添字をもつものを\textbf{\bookterm{double}{二重数列}}という.
一変数で使った階差と帰納法を2方向へ広げた考え方である.

\begin{bookpoint}
  \textbf{二つの添字と初期条件}\par
  \bookterm{initial}{初期条件}となる境界と,項を求める順序を確かめる.
  二方向の規則を別々に与えるなら,同じ位置を二通りに計算して一致する必要がある.
\end{bookpoint}
格子経路と二項係数の資料は参考文献のNIST DLMF §26.3へ.
